\documentclass[11pt, a4paper]{article}

\usepackage[utf8]{inputenc}
\usepackage[T1]{fontenc}
\usepackage[spanish]{babel}
\usepackage{amsmath, amssymb, amsthm}
\usepackage{geometry}
\usepackage{hyperref}
\usepackage{booktabs}
\usepackage{tabularx}
\usepackage{float}
\usepackage{xcolor}

\geometry{margin=2.5cm}

\hypersetup{
  colorlinks=true,
  linkcolor=blue!60!black,
  citecolor=blue!60!black,
  urlcolor=blue!60!black,
}

\theoremstyle{definition}
\newtheorem{definition}{Definición}[section]
\newtheorem{remark}{Observación}[section]
\newtheorem{theorem}{Teorema}[section]

\newcommand{\R}{\mathbb{R}}
\newcommand{\E}{\mathbb{E}}
\newcommand{\Prob}{\mathbb{P}}
\newcommand{\norm}[1]{\left\|#1\right\|}
\newcommand{\abs}[1]{\left|#1\right|}
\newcommand{\inner}[2]{\langle #1,\, #2 \rangle}
\newcommand{\eps}{\varepsilon}

\title{\textbf{Arquitectura matemática del código}\\[0.5em]
\large Aplicación del control óptimo de campo medio a clasificación binaria\\
\normalsize Basado en Daudin \& Delarue, \textit{arXiv:2507.08486}}
\date{}

% ============================================================
\date{}
\begin{document}
\maketitle
\tableofcontents
\newpage

% ============================================================
\section{Datos: generación y preprocesado}
\label{sec:datos}

\subsection{Generación del conjunto \texttt{make\_moons}}

Se generan $N$ pares $(x_i, y_i)$ con la función \texttt{make\_moons} de
\texttt{scikit-learn}. Cada punto pertenece a una de dos clases entrelazadas:
$y_i \in \{0, 1\}$, con $N/2$ puntos por clase.

En primer lugar, se generan coordenadas limpias $x_i \in \R^2$: para $\theta \in [0, \pi]$, 
$x = (\cos \theta, \sin \theta)$ pertenece a la clase $0$, y $x = (1 - \cos \theta, 1 - \sin \theta - 0.5)$ 
pertenece a la clase $1$. Esto produce dos lunas entrelazadas.

A las coordenadas limpias se les añade ruido gaussiano independiente en cada
dimensión:
\begin{equation}
  x_i \;\leftarrow\; x_i + \eta_i, \qquad
  \eta_i \sim \mathcal{N}(0,\, \sigma^2 I_2), \quad \sigma = 0.12.
\end{equation}

\subsection{Normalización (\texttt{StandardScaler})}

Tras añadir el ruido, se aplica normalización estándar componente a componente.
Sea $\mu_k$ y $s_k$ la media y desviación típica de la componente $k$ sobre los
$N$ puntos. La transformación es:
\begin{equation}
  \tilde{x}_i^{(k)} = \frac{x_i^{(k)} - \mu_k}{s_k}, \qquad k = 1, 2.
\end{equation}
Tras la normalización, cada coordenada tiene media $0$ y varianza $1$. Esto
estabiliza el entrenamiento al colocar los datos en una escala homogénea.

En adelante se denota $x_0^{i} \in \R^2$ al punto normalizado $i$-ésimo y
$\gamma_0 = \frac{1}{N}\sum_{i=1}^N \delta_{(x_0^{i}, y_0^{i})}$ a la medida
empírica de los datos de entrenamiento. La base de datos es $\{(x_0^{i}, y_0^{i})\}_{i=1}^N$.

% ============================================================
\section{Arquitectura del modelo}
\label{sec:arquitectura}

El modelo es una \emph{Mean-Field ResNet} compuesta de dos bloques:
\begin{enumerate}
  \item Un \textbf{campo vectorial de campo medio} (\texttt{MeanFieldVelocity})
        que define la dinámica de una ODE,
  \item un \textbf{clasificador lineal} que opera sobre el estado final de la ODE.
\end{enumerate}

\subsection{Campo vectorial prototípico}
\label{sec:velocity}

Siguiendo el Ejemplo~1.1 del paper \cite{daudin2025}, se utiliza el campo
vectorial prototípico
\begin{equation}
  b(x, a) = \sigma(a_1 \cdot x + a_2)\, a_0,
  \qquad
  \sigma = \tanh,
\end{equation}
donde el \emph{parámetro de control} de cada neurona $m$ es
\begin{equation}
  a^m = (a_0^m,\, a_1^m,\, a_2^m,\, b_1^m)
  \;\in\; \R^2 \times \R^2 \times \R \times \R \;=\; \R^6.
\end{equation}
Los roles de cada componente son:
\begin{itemize}
  \item $a_1^m \in \R^2$: \emph{pesos de entrada}; definen la dirección del
        espacio de features que la neurona $m$ ``proyecta''. La
        pre-activación es $a_1^m \cdot x \in \R$.
  \item $a_2^m \in \R$ (columna temporal de $W_1$): coeficiente que escala la
        contribución explícita del tiempo $t \in [0,T]$ a la pre-activación.
  \item $b_1^m \in \R$: sesgo fijo de la neurona $m$ en $t=0$.
  \item $a_0^m \in \R^2$: \emph{pesos de salida}; escalan la contribución de la
        neurona al campo vectorial en $\R^2$. Su norma $\norm{a_0^m}_2$ mide la
        \emph{importancia} de la neurona.
\end{itemize}

\subsection{Campo efectivo de campo medio}
\label{sec:campo_efectivo}

En la aproximación de $M$ partículas (neuronas), la medida de control
$\nu_t \approx \frac{1}{M}\sum_{m=1}^M \delta_{a^m}$ produce el campo
efectivo:
\begin{equation}
  \label{eq:campo}
  F(x, t) = \frac{1}{M}\sum_{m=1}^{M} b(x, a^m(t))
  = \frac{1}{M}\sum_{m=1}^{M}
    \sigma\!\bigl(a_1^m \cdot x + a_2^m\, t + b_1^m\bigr)\, a_0^m.
\end{equation}

\paragraph{Diferencia con el paper.}
El paper optimiza sobre trayectorias arbitrarias $(a^m(t))_{t \in [0,T]}$.
En la implementación, los parámetros $(a_1^m, a_0^m, b_1^m)$ son
\textbf{estáticos} (independientes de $t$), y la dependencia temporal se
introduce \emph{explícitamente} añadiendo $t$ como entrada adicional:
\begin{equation}
  \text{pre-activación}^m(x,t) = a_1^m \cdot x + a_2^m \cdot t + b_1^m.
\end{equation}
Esto restringe la familia de controles a aquellos cuya trayectoria
$a_2^m \mapsto a_2^m \cdot t + b_1^m$ varía \emph{linealmente} en $t$,
pero preserva la dependencia temporal del campo $F(x,t)$.

\paragraph{Implementación matricial.}
Con $M=64$ neuronas y $d_1=2$:
\begin{align}
  W_1 &\in \R^{M \times (d_1+1)}, \quad
    [W_1]_{m,:} = (a_1^m,\, a_2^m), \quad
    b_1 \in \R^M, \\
  W_0 &\in \R^{d_1 \times M}, \quad
    [W_0]_{:,m} = a_0^m \quad \text{(sin bias)}.
\end{align}
Para $N$ puntos simultáneos $X \in \R^{N \times d_1}$:
\begin{equation}
  h = \tanh\!\bigl([\,X,\; t\,\mathbf{1}\,] \, W_1^\top + b_1\bigr)
  \;\in\; \R^{N \times M},
  \qquad
  F(X, t) = \frac{1}{M}\, h\, W_0^\top \;\in\; \R^{N \times d_1}.
\end{equation}
El factor $1/M$ se absorbe en la escala de inicialización de $W_0$.

\paragraph{Total de parámetros de la ODE.}
\begin{equation}
  \underbrace{M \cdot (d_1+1)}_{\text{pesos } W_1} +
  \underbrace{M}_{\text{bias } b_1} +
  \underbrace{d_1 \cdot M}_{\text{pesos } W_0}
  = 64 \cdot 3 + 64 + 2 \cdot 64 = 384.
\end{equation}

\subsection{ODE de transporte y flujo}
\label{sec:ode}

El estado $X_t \in \R^2$ de cada punto de dato evoluciona según la ODE de
primer orden:
\begin{equation}
  \label{eq:ode}
  \frac{dX_t}{dt} = F(X_t, t), \qquad X_0 = x_0, \qquad t \in [0, T],
  \quad T = 1.
\end{equation}
La solución define el \emph{flujo} $\phi_t : \R^2 \to \R^2$ con
$X_t = \phi_t(x_0)$. La distribución de features evoluciona como la
\emph{medida push-forward}:
\begin{equation}
  \gamma_t = (\phi_t)_\# \gamma_0,
\end{equation}
lo que es equivalente a la ecuación de continuidad (ec.~1.3 del paper):
\begin{equation}
  \partial_t \gamma_t + \operatorname{div}_x\!\bigl(F(x,t)\,\gamma_t\bigr) = 0.
\end{equation}
Esta PDE garantiza conservación de masa: los puntos no se crean ni se destruyen,
sino que son \emph{transportados} por el campo $F$.

\subsection{Integración numérica: Runge-Kutta de orden 4}
\label{sec:rk4}

La ODE~\eqref{eq:ode} se integra con $n_\text{steps} = 10$ pasos de
Runge-Kutta 4 (RK4), con paso $\Delta t = T / n_\text{steps} = 0.1$.
En cada paso $k$ (con $t_k = k \cdot \Delta t$):
\begin{align}
  k_1 &= F(X_{t_k},\; t_k), \\
  k_2 &= F\!\left(X_{t_k} + \tfrac{\Delta t}{2}\,k_1,\;
                  t_k + \tfrac{\Delta t}{2}\right), \\
  k_3 &= F\!\left(X_{t_k} + \tfrac{\Delta t}{2}\,k_2,\;
                  t_k + \tfrac{\Delta t}{2}\right), \\
  k_4 &= F\!\left(X_{t_k} + \Delta t\,k_3,\;
                  t_k + \Delta t\right), \\
  X_{t_{k+1}} &= X_{t_k} + \frac{\Delta t}{6}(k_1 + 2k_2 + 2k_3 + k_4).
\end{align}

El error global de RK4 es $\mathcal{O}(\Delta t^4) = \mathcal{O}(10^{-4})$,
cuatro órdenes de magnitud mejor que Euler ($\mathcal{O}(\Delta t) = 0.1$),
al coste de cuatro evaluaciones de $F$ por paso.

\subsection{Clasificador lineal}
\label{sec:clasificador}

Tras integrar la ODE hasta $t = T$, se obtiene el estado final
$X_T \in \R^{N \times 2}$. El clasificador lineal calcula el logit:
\begin{equation}
  \ell_i = W \cdot X_T^{i} + b,
  \qquad W \in \R^{1 \times 2},\; b \in \R,
\end{equation}
y la probabilidad predicha es:
\begin{equation}
  \hat{p}_i = \Prob(y_i = 1 \mid X_T^{i}) = \sigma(\ell_i)
  = \frac{1}{1 + e^{-\ell_i}}.
\end{equation}

\paragraph{Parámetros del clasificador.}
\begin{equation}
  \underbrace{1 \times 2}_{\text{pesos } W} + \underbrace{1}_{\text{bias } b}
  = 3 \text{ parámetros},
\end{equation}
lo que representa menos del $1\%$ del total de parámetros del modelo ($387$).

\paragraph{Intuición.}
La ODE aprende a \emph{transportar} la distribución $\gamma_0$ (lunas
entrelazadas, no separables linealmente) hasta $\gamma_T$ (clases separadas).
El clasificador lineal simplemente lee el resultado: toda la complejidad
no lineal está en el flujo.

% ============================================================
\section{Función de pérdida}
\label{sec:perdida}

\subsection{Coste total}

La función de pérdida que se minimiza es:
\begin{equation}
  \label{eq:loss}
  J(\theta) = \underbrace{\mathcal{L}_{\text{BCE}}(\theta)}_{\text{coste terminal}}
            + \varepsilon \cdot \underbrace{\mathcal{E}(\theta)}_{\text{penalización entrópica}},
\end{equation}
donde $\theta$ denota el conjunto completo de parámetros
$(W_1, b_1, W_0, W, b)$ y $\varepsilon > 0$ es el coeficiente de
regularización.

\subsection{Entropía cruzada binaria (BCE)}
\label{sec:bce}

El coste terminal es la entropía cruzada binaria (Binary Cross-Entropy) con
logits:
\begin{equation}
  \label{eq:bce}
  \mathcal{L}_{\text{BCE}}(\theta)
  = -\frac{1}{N}\sum_{i=1}^{N}
    \Bigl[ y_i \log \sigma(\ell_i) + (1 - y_i)\log(1 - \sigma(\ell_i)) \Bigr],
\end{equation}
donde $\ell_i = W \cdot X_T^{i} + b$ y $\sigma(\ell_i) = (1 + e^{-\ell_i})^{-1}$.
En términos del paper \cite{daudin2025}, este es el coste terminal
$L(X_T, Y)$ del problema de control óptimo:
\begin{equation}
  \min_{\nu_t} \E\bigl[L(X_T, Y)\bigr] + 
  \eps \int_0^T \cdot E(\nu_t \mid \nu^\infty)\,dt,
\end{equation}
donde $E(\nu_t \mid \nu^\infty)$ es la divergencia KL relativa al prior
$\nu^\infty$.

\subsection{Penalización entrópica}
\label{sec:entropica}

\paragraph{Motivación teórica.}
El paper regula la medida de control $\nu_t$ mediante la divergencia KL
respecto al prior de Gibbs:
\begin{equation}
  \nu^\infty \propto \exp(-\ell(a)), \qquad
  \ell(a) = c_1 \abs{a}^4 + c_2 \abs{a}^2,
  \quad c_1 = 0.05,\; c_2 = 0.5.
\end{equation}
El término cuártico $c_1 \abs{a}^4$ es esencial: satisface la condición de
\emph{supercoercividad} ($\nabla^2 \ell(a) \geq c(1 + \abs{a}^2)I$) que
garantiza la desigualdad log-Sobolev para $\nu^\infty$ y por tanto la
condición Polyak-Łojasiewicz (véase Sección~\ref{sec:pl}).
Un potencial puramente cuadrático ($c_1 = 0$) no sería suficiente.

La KL completa es:
\begin{equation}
  \operatorname{KL}(\nu_t \Vert \nu^\infty)
  = \underbrace{\E_{\nu_t}[\ell(a)]}_{\text{término de energía}}
  - \underbrace{H(\nu_t)}_{\text{entropía diferencial}}.
\end{equation}

\paragraph{Limitación de la implementación.}
En la implementación, $\nu_t = \frac{1}{M}\sum_{m=1}^M \delta_{\theta^m}$
es una suma de deltas de Dirac sobre los parámetros puntuales. La entropía
diferencial de una medida discreta respecto a un prior continuo es
$-\infty$, por lo que la KL completa no es computable.

Sólo es accesible el \textbf{término de energía}:
\begin{equation}
  \label{eq:energia}
  \mathcal{E}(\theta)
  = \frac{1}{N_p}\sum_{j=1}^{N_p}
    \Bigl[c_1\, \theta_j^4 + c_2\, \theta_j^2\Bigr],
  \quad c_1 = 0.05,\; c_2 = 0.5,
\end{equation}
donde la suma recorre los $N_p = 384$ parámetros de \texttt{MeanFieldVelocity}
(se \emph{excluyen} los 3 parámetros del clasificador, que son parte del
coste terminal, no del control). Dividir por $N_p$ hace que el valor no
dependa del tamaño del modelo.

\begin{remark}
  En términos prácticos, $\mathcal{E}(\theta)$ es equivalente a regularización
  $L^4 + L^2$ (\emph{weight decay} polinomial). Este término proporciona el
  \textbf{prior energético} $\nu^\infty$, pero no el término de entropía
  $-H(\nu_t)$ de la KL completa. Para los experimentos B, E y F, el término
  de entropía se implementa mediante \textbf{SGLD} (ruido de Langevin en el
  gradiente, véase Sección~\ref{sec:sgld}).
\end{remark}

\paragraph{Estructura de $a^m$ para la penalización.}
Recordando que cada neurona $m$ tiene parámetro
$a^m = (a_0^m, a_1^m, a_2^m, b_1^m) \in \R^6$, la norma que aparece en
$\ell(a)$ es:
\begin{equation}
  \abs{a^m}^2 = \norm{a_0^m}^2 + \norm{a_1^m}^2 + (a_2^m)^2 + (b_1^m)^2.
\end{equation}
El término de energía es la esperanza de $\ell$ bajo la medida empírica
$\nu_t = \frac{1}{M}\sum_{m=1}^M \delta_{a^m}$:
\begin{equation*}
  \mathbb{E}_{\nu_t}[\ell(a)]
  = \frac{1}{M}\sum_{m=1}^M \ell(a^m)
  = \frac{1}{M}\sum_{m=1}^M \bigl(c_1|a^m|^4 + c_2|a^m|^2\bigr).
\end{equation*}
Como $|a^m|^2 = \sum_j (a^m_j)^2$ descompone en escalares, esta expresión
es equivalente a aplicar $c_1\theta_j^4 + c_2\theta_j^2$ a cada parámetro
escalar $\theta_j$ y promediar — de ahí la forma de la ecuación~\eqref{eq:energia}.

\begin{remark}[Discrepancia en el término cuártico]
  El paper usa $|a|^4 = \bigl(\sum_j\theta_j^2\bigr)^2$; la implementación
  calcula $\sum_j\theta_j^4$, que omite los términos cruzados $\theta_i^2\theta_j^2$.
  Ambos son supercoercivos, por lo que las garantías teóricas se preservan.
\end{remark}

% ============================================================
\section{Entrenamiento}
\label{sec:entrenamiento}

La función \texttt{train()} implementa tres modos de optimización, cada uno
adecuado a un experimento distinto:

\begin{center}
\begin{tabular}{llll}
\toprule
Modo & Parámetro & Experimentos & Objetivo \\
\midrule
Adam + cosine annealing & (defecto) & A, D & Convergencia rápida \\
SGD + lr constante & \texttt{use\_sgd=True} & C & Verificación PL sin artefactos \\
pSGLD + cosine annealing & \texttt{use\_sgld=True} & B, E, F, G & Muestreo de $\nu_t^*$ \\
\bottomrule
\end{tabular}
\end{center}

\subsection{Optimizador Adam}
\label{sec:adam}

Se utiliza el optimizador Adam \cite{kingma2015adam} con tasa de aprendizaje
$\alpha_0 = 0.01$, $\beta_1 = 0.9$, $\beta_2 = 0.999$, $\eps_{\text{Adam}} = 10^{-8}$.
En cada paso de gradiente $s$, para cada parámetro escalar $\theta$:
\begin{align}
  g_s &= \nabla_\theta J(\theta^{s-1}), \\
  m_s &= \beta_1 m_{s-1} + (1 - \beta_1)\, g_s
        \qquad \text{(primer momento, media exponencial)}, \\
  v_s &= \beta_2 v_{s-1} + (1 - \beta_2)\, g_s^2
        \qquad \text{(segundo momento, varianza exponencial)}, \\
  \hat{m}_s &= \frac{m_s}{1 - \beta_1^s}
        \qquad \text{(corrección de sesgo)}, \\
  \hat{v}_s &= \frac{v_s}{1 - \beta_2^s}, \\
  \theta^s &= \theta^{s-1} - \alpha_s\, \frac{\hat{m}_s}{\sqrt{\hat{v}_s} + \eps_{\text{Adam}}}.
\end{align}
La corrección de sesgo es importante en los primeros pasos, cuando
$m_s \approx (1-\beta_1)g_1$ y $v_s \approx (1-\beta_2)g_1^2$ son aún
pequeños. Sin corrección, el paso efectivo sería demasiado grande.

\subsection{SGD con tasa de aprendizaje constante}
\label{sec:sgd}

Para el Experimento~C (verificación de la condición PL), se usa SGD puro con
tasa de aprendizaje \emph{constante} $\alpha = 0.01$ y sin ruido:
\begin{equation}
  \theta^s = \theta^{s-1} - \alpha\, \nabla_\theta J(\theta^{s-1}).
\end{equation}
La motivación es que la condición PL es una propiedad \textbf{geométrica} del
paisaje de $J$, independiente del optimizador. Con cosine annealing,
$\alpha_s \to 0$ al final aplana $J(\theta^s) - J^*$ artificialmente, lo que
contamina la estimación de $\hat{\mu}$. Con lr constante, $J^*$ es el mínimo
geométrico real del gradient flow.

\subsection{SGLD precondicionado (pSGLD)}
\label{sec:sgld}

Para los Experimentos B, E y F, se usa \textbf{pSGLD} (\emph{preconditioned
Stochastic Gradient Langevin Dynamics}, Li et al.~\cite{li2016psgld}):
una extensión de SGLD que utiliza \textbf{Adam} como optimizador base y
acopla el ruido de Langevin al precondicionador de Adam.

\paragraph{Motivación: por qué Adam como base, no SGD.}
En este problema los gradientes de la BCE son del orden de $10^{-5}$.
SGD con $\eta = 0.01$ produce pasos de tamaño $\sim 10^{-7}$, insuficientes
para converger. Adam adapta la tasa de aprendizaje por componente mediante su
precondicionador, garantizando convergencia donde SGD falla.

\paragraph{Fórmula de actualización pSGLD.}
En cada paso $s$, tras el paso de Adam:
\begin{equation}
  \label{eq:psgld}
  \theta^s \leftarrow \underbrace{\text{Adam}(\theta^{s-1}, \nabla J)}_{\text{paso determinístico}}
  \;+\; \underbrace{\sqrt{2\,\eta_s\,\varepsilon \cdot M_s}\;\xi^s}_{\text{ruido precondicionado}},
  \qquad \xi^s \sim \mathcal{N}(0, I),
\end{equation}
donde $M_s \in \R^{N_p}$ es el vector de precondicionadores componente a
componente:
\begin{equation}
  \label{eq:precond}
  M_s[j] = \min\!\left(\frac{1}{\sqrt{\hat{v}_s[j]} + \delta},\; 1\right),
  \qquad
  \hat{v}_s[j] = \frac{v_s[j]}{1 - \beta_2^s},
\end{equation}
con $v_s[j]$ el segundo momento exponencial de Adam para el parámetro $j$,
$\beta_2 = 0.999$ y $\delta = 10^{-8}$.

\paragraph{Corrección de sesgo.}
El término $\hat{v}_s[j] = v_s[j]/(1-\beta_2^s)$ es el segundo momento
sesgado-corregido: en los primeros pasos, $v_s[j] \approx (1-\beta_2)g_1^2$
es pequeño aunque el gradiente no lo sea; sin corrección, $M_s[j]$ sería
artificialmente grande y el ruido explotaría. La corrección garantiza que
$\hat{v}_s[j] \to \mathbb{E}[g^2]$ a medida que $s$ crece.

\paragraph{Cota superior en $M_s$ y justificación.}
Sin la cota $M_s \leq 1$, los parámetros con dirección plana
($\hat{v}_s[j] \approx 0$) producirían $M_s[j] \approx 1/\delta = 10^8$,
y el ruido $\sqrt{2\eta\varepsilon \cdot 10^8}$ explotaría incluso para
$\varepsilon$ pequeño. La cota $M_s \leq 1$ corresponde a usar un
hiperparámetro de estabilización $\lambda \geq 1$ (distinto del $\delta$
de Adam) en la fórmula $1/(\sqrt{\hat{v}} + \lambda)$, como sugieren
Li et al.\ en su implementación original. Con esta cota, el ruido
satisface $\|\text{ruido}\| \leq \sqrt{2\eta_s\varepsilon}$: nunca supera
el nivel isotrópico de SGLD estándar.

\paragraph{Necesidad del acoplamiento ruido--precondicionador.}
Inyectar ruido isotrópico ($\sigma$ constante para todos los $j$) después
del paso de Adam desacopla la parte estocástica del precondicionador y
destruye la distribución estacionaria: la cadena de Markov ya \emph{no}
converge a $\nu^* \propto \exp(-J(\theta)/\varepsilon)$. Con el acoplamiento
de~\eqref{eq:psgld}--\eqref{eq:precond}, la SDE efectiva es la Langevin
precondicionada
\begin{equation}
  d\theta = -\nabla J\,dt + \sqrt{2\varepsilon M}\,dW,
\end{equation}
cuya medida estacionaria es $\nu^* \propto \exp(-J(\theta)/\varepsilon)$
bajo el precondicionador $M$ (resultado estándar de difusiones precondicionadas).

\paragraph{Resultado.}
El ruido precondicionado implementa el \textbf{término de entropía}
$-H(\nu_t)$ de la KL, haciendo que los parámetros exploren la distribución
de Gibbs en lugar de colapsar al MAP. El término energético $\mathcal{E}$
en~\eqref{eq:energia} sigue actuando como el prior $\nu^\infty$.
Con $\varepsilon = 0$ el ruido es exactamente cero (Adam estándar).

La tasa de aprendizaje sigue cosine annealing ($\eta_s \to 0$), satisfaciendo
las condiciones de Robbins-Monro: $\sum_s \eta_s = \infty$,
$\sum_s \eta_s^2 < \infty$.

\subsection{Cosine Annealing}
\label{sec:cosine}

La tasa de aprendizaje $\alpha_s$ no es constante: sigue un calendario de
\emph{cosine annealing} \cite{loshchilov2017sgdr} con $\eta_{\min} = 0$ y
$T_{\max} = n_\text{epochs}$:
\begin{equation}
  \alpha_s = \frac{\alpha_0}{2}\Bigl(1 + \cos\!\Bigl(\frac{\pi s}{T_{\max}}\Bigr)\Bigr),
  \qquad s = 0, 1, \ldots, T_{\max}.
\end{equation}
La tasa arranca en $\alpha_0 = 0.01$, decrece suavemente hasta $0$ al final
del entrenamiento. En Adam (experimentos A, D) y SGLD (B, E, F), el
aplanamiento de las curvas de pérdida en las últimas épocas se debe al
scheduler, no a una violación de la condición PL. Por eso el Experimento~C
usa SGD con lr constante (Sección~\ref{sec:sgd}).

\subsection{Gradient Clipping}
\label{sec:clipping}

Para evitar gradientes explosivos en las primeras épocas (cuando los
parámetros son aún aleatorios y el campo $F$ puede ser grande), se aplica
\emph{gradient clipping} por norma global:
\begin{equation}
  g \;\leftarrow\;
  \begin{cases}
    g & \text{si } \norm{g} \leq G_{\max}, \\
    G_{\max}\, \dfrac{g}{\norm{g}} & \text{si } \norm{g} > G_{\max},
  \end{cases}
  \qquad G_{\max} = 5.
\end{equation}
Aquí $g = \nabla_\theta J(\theta)$ es el vector de gradientes concatenado
sobre todos los parámetros, y la norma es la euclídea global. El clipping
escala el vector entero (no componente a componente), preservando la
dirección del gradiente.

En Adam, el segundo momento $\hat{v}_s$ hace de divisor adaptativo, pero sin
clipping la actualización puede ser muy grande en la primera época cuando
$\hat{v}_s$ todavía no ha acumulado información.

\subsection{Retropropagación}
\label{sec:backprop}

El gradiente $\nabla_\theta J$ se calcula automáticamente con
\texttt{loss.backward()} de PyTorch. El grafo computacional incluye:
\begin{enumerate}
  \item Diferenciación a través del clasificador lineal ($\partial J / \partial W$,
        $\partial J / \partial b$).
  \item Diferenciación a través del RK4 (cuatro evaluaciones de $F$ por paso,
        $n_\text{steps} = 10$ pasos): $\partial J / \partial X_T \cdot
        \partial X_T / \partial \theta$ mediante la regla de la cadena.
  \item Diferenciación de la penalización entrópica
        $\partial \mathcal{E} / \partial \theta_j = 4c_1 \theta_j^3 + 2c_2 \theta_j$.
\end{enumerate}
PyTorch construye el grafo dinámicamente en el paso \emph{forward}, por lo que
no es necesario implementar el adjunto del integrador manualmente.

\subsection{Hiperparámetros de entrenamiento}

\begin{table}[H]
\centering
\begin{tabularx}{\textwidth}{l l l X}
\toprule
Hiperparámetro & Valor & Experimentos & Descripción \\
\midrule
$N$ & 400 & todos & Puntos de entrenamiento \\
$d_1$ & 2 & todos & Dimensión del espacio de features \\
$M$ & 64 & todos & Neuronas / partículas \\
$T$ & 1.0 & todos & Horizonte temporal de la ODE \\
$n_\text{steps}$ & 10 & todos & Pasos RK4 \\
$\varepsilon$ & 0.01 & todos & Temperatura entrópica / regularización \\
$\alpha_0$ & 0.01 & todos & Tasa de aprendizaje inicial \\
$\beta_1, \beta_2$ & 0.9, 0.999 & A, D & Momentos de Adam \\
$n_\text{epochs}$ & 800 (A), 700 (B–G) & todos & Épocas de entrenamiento \\
$G_{\max}$ & 5.0 & todos & Umbral de gradient clipping \\
$c_1$ & 0.05 & todos & Coeficiente cuártico del prior \\
$c_2$ & 0.5 & todos & Coeficiente cuadrático del prior \\
\midrule
\multicolumn{4}{l}{\textit{Solo pSGLD (experimentos B, E, F, G):}} \\
$\sqrt{2\eta_s\varepsilon M_s}$ & variable & B, E, F, G & Escala del ruido precondicionado (por componente) \\
$\lambda$ & $1$ & B, E, F, G & Cota superior de $M_s$ (evita explosión en dirs.\ planas) \\
\bottomrule
\end{tabularx}
\caption{Hiperparámetros del modelo y el entrenamiento.}
\end{table}

% ============================================================
\section{Fundamentos teóricos}
\label{sec:teoria}

\subsection{Condición Polyak-Łojasiewicz (PL)}
\label{sec:pl}

El paper \cite{daudin2025} demuestra (Meta-Teorema~2) que para $\varepsilon > 0$
y condiciones iniciales en un conjunto abierto y denso $\mathcal{O}$, la
función de pérdida satisface la \emph{desigualdad de Polyak-Łojasiewicz}
(PL) local:
\begin{equation}
  \label{eq:pl}
  \norm{\nabla J(\theta)}^2 \;\geq\; 2\mu \cdot (J(\theta) - J^*),
  \qquad \mu > 0,
\end{equation}
donde $J^* = \inf_\theta J(\theta)$ es el valor óptimo. Esta condición
implica convergencia \emph{exponencial} del gradient descent con tasa
de aprendizaje $\eta$:
\begin{equation}
  J(\theta^s) - J^* \;\leq\; (1 - 2\eta\mu)^s \cdot (J(\theta^0) - J^*).
\end{equation}

\begin{remark}
  La condición PL es más débil que la convexidad fuerte: no implica que el
  espacio de parámetros sea convexo ni que el mínimo sea único en sentido
  geométrico, pero sí garantiza que cualquier punto estacionario es global.
  La clave es que $\varepsilon > 0$: con $\varepsilon = 0$ no hay garantía.
\end{remark}

\subsection{Genericidad del minimizador}
\label{sec:genericidad}

El Meta-Teorema~1 del paper establece que para un conjunto abierto y denso
$\mathcal{O}$ de condiciones iniciales $\gamma_0$, el problema de control
tiene un único minimizador \emph{estable}. Esto predice robustez empírica:
distintas inicializaciones de parámetros y distintos datasets deberían
producir el mismo $J^*$ (a igualdad de clase de problema).

% ============================================================
\section{Descripción matemática de los experimentos}
\label{sec:experimentos}

\subsection{Experimento A: Evolución de $\gamma_t$}

\textbf{Objetivo.} Visualizar que la ODE transforma $\gamma_0$ (lunas no
separables) en $\gamma_T$ (clases separadas).

\textbf{Protocolo.} Se entrena un modelo con $\varepsilon = 0.01$, $M = 64$,
$T = 1$, $n_\text{epochs} = 800$. Tras el entrenamiento, se integra la ODE
guardando $X_t$ en $t \in \{0, T/4, T/2, 3T/4, T\}$. Se visualiza la
distribución $\gamma_t$ en cada instante.

\textbf{Magnitudes visualizadas.}
\begin{itemize}
  \item Snapshots $\gamma_t$ para $t \in \{0, 0.25, 0.5, 0.75, 1.0\}$.
  \item Trayectorias $t \mapsto X_t^{(i)}$ para 40 partículas seleccionadas.
  \item Curvas de convergencia: $J$, BCE, $\mathcal{E}$ vs. época.
  \item Frontera de decisión final en $\R^2$.
\end{itemize}

\textbf{Resultado esperado.} $\gamma_T$ debe ser linealmente separable con
accuracy $\approx 1.0$.

\subsection{Experimento B: Efecto del parámetro $\varepsilon$}

\textbf{Objetivo.} Estudiar el efecto del coeficiente de regularización
entrópica sobre la convergencia, las fronteras de decisión y el campo de
velocidad.

\textbf{Protocolo.} Se entrena un modelo para cada
$\varepsilon \in \{0, 10^{-3}, 10^{-2}, 0.1, 0.5\}$ con \textbf{pSGLD}
(\texttt{use\_sgld=True}, Li et al.~\cite{li2016psgld}), manteniendo la
misma semilla de inicialización. El uso de pSGLD es coherente con el rol
de $\varepsilon$ como temperatura entrópica: el ruido precondicionado
$\sqrt{2\eta_s\varepsilon M_s}\,\xi$ hace que los parámetros exploren
$\nu^* \propto \exp(-J(\theta)/\varepsilon)$. Con $\varepsilon = 0$ el
ruido es exactamente cero (Adam estándar).

\textbf{Magnitudes analizadas.}

\begin{itemize}
  \item \textbf{B1 — Curvas de convergencia:} $J$ total, accuracy y
        $\mathcal{E}$ vs.\ época para cada $\varepsilon$. Los ejes de $J$
        y $\mathcal{E}$ están recortados ($J \leq 0.5$,
        $\mathcal{E} \leq 0.15$) para que $\varepsilon=0.5$ no domine la
        escala; las curvas de accuracy muestran la señal bruta (trazo fino)
        y una media móvil de ventana 15 (trazo grueso).
        Predicción: $J^*$ crece con $\varepsilon$ (la regularización penaliza
        la solución óptima ``pura''), pero todas las $\varepsilon$ convergen a
        accuracy $\approx 1.0$.
  \item \textbf{B2 — Fronteras de decisión y campo de velocidad:} la fila
        superior muestra la región de decisión $\{x : \hat{p}(x) \geq 0.5\}$
        para cada $\varepsilon$; la fila inferior muestra el quiver plot del
        campo $F(x, t=0.5)$ normalizado en dirección (color = magnitud).
        Predicción: fronteras más suaves para $\varepsilon$ mayor;
        campo más regular.
\end{itemize}

\textbf{Trade-off empírico $\varepsilon$ vs. $\hat{\mu}_{\text{PL}}$.}
Empíricamente se observa que $\hat{\mu}_{\text{PL}}$ no crece
monotónicamente con $\varepsilon$. Esto no contradice el paper: la teoría
garantiza $\mu > 0$ para todo $\varepsilon > 0$, pero no que $\mu$ sea
monótono. Con $\varepsilon$ grande, $J^*$ es mayor (más penalización) y el
exceso de coste $(J - J^*)$ en el denominador de la ratio PL puede crecer
más que $\norm{\nabla J}^2$, reduciendo $\hat{\mu}$.

\subsection{Experimento C: Verificación de la condición PL}

\textbf{Objetivo.} Verificar empíricamente la desigualdad~\eqref{eq:pl}.

\textbf{Protocolo.} Experimento \textbf{autocontenido}: entrena sus propios
modelos con \textbf{SGD + lr constante} (\texttt{use\_sgd=True},
$\varepsilon \in \{0, 10^{-3}, 10^{-2}, 0.1, 0.5\}$, 700 épocas). No
reutiliza modelos del Experimento~B. La motivación es que la condición PL es
una propiedad \textbf{geométrica} del paisaje de $J$, independiente del
optimizador: con cosine annealing, $\alpha_s \to 0$ aplana $\delta_s =
J(\theta^s) - J^*$ artificialmente y contamina $\hat{\mu}$. Con lr constante,
$J^*$ es el mínimo geométrico genuino del gradient flow.

En cada época $s$, se registran:
\begin{equation}
  \text{excess cost}: \quad \delta_s = J(\theta^s) - J^*,
  \qquad
  \text{grad norm}: \quad G_s = \norm{\nabla J(\theta^s)}^2.
\end{equation}
El estimador conservador de $\mu$ se define como el percentil 10 del ratio
\begin{equation}
  \hat{\mu}_{\text{PL}} = \operatorname{P}_{10}\!\left\{
    \frac{G_s}{2\,\delta_s} : s = 1, \ldots, n_\text{epochs}
  \right\}.
\end{equation}

\textbf{Paneles generados.}
\begin{itemize}
  \item \textbf{C1 — Log-log:} $G_s$ vs.\ $\delta_s$. Los puntos deben quedar
        por encima de la recta $G = 2\hat{\mu}\,\delta$.
  \item \textbf{C2 — Semilog:} $\delta_s$ vs.\ época. Caída exponencial
        confirma la convergencia predicha por~\eqref{eq:pl}.
  \item \textbf{C3 — Barras:} $\hat{\mu}_{\text{PL}}$ para cada $\varepsilon$.
        Confirma $\hat{\mu} > 0$ para todo $\varepsilon > 0$.
  \item \textbf{C4 — Ratio PL:} $G_s / (2\delta_s)$ vs.\ época (escala log),
        que debe mantenerse $\geq \hat{\mu}$ en todo momento.
\end{itemize}

\subsection{Experimento D: Genericidad del minimizador}

\textbf{Objetivo.} Verificar empíricamente el Meta-Teorema~1: que el
minimizador es robusto a cambios en la inicialización y en el dataset.

\textbf{Protocolo.} Se ejecutan tres sub-experimentos con
$n_\text{seeds} = 10$, 500 épocas, Adam y $\varepsilon \in \{0, 0.01\}$:
\begin{itemize}
  \item \textbf{D1 — $\gamma_0$ fija, inicialización variable:}
        \texttt{data\_seed}$=42$, \texttt{init\_seed}$\in\{0,\ldots,9\}$.
  \item \textbf{D2 — Inicialización fija, $\gamma_0$ variable:}
        \texttt{init\_seed}$=2$, \texttt{data\_seed}$\in\{0,\ldots,9\}$.
  \item \textbf{D3 — Ambas aleatorias:}
        \texttt{data\_seed}$=$\texttt{init\_seed}$=s$.
\end{itemize}

\textbf{Resultado esperado.} D3 tiene la banda de pérdida más ancha (combina
ambas fuentes de variabilidad). Con $\varepsilon = 0.01$, la condición PL
garantiza $\hat{\mu} > 0$ para todo run, pero la regularización no reduce
necesariamente la varianza entre seeds (el paper garantiza convergencia desde
cualquier punto, no menor dispersión). Las fronteras de decisión de D2 y D3
son topológicamente similares entre sí, verificando el Meta-Teorema~1.
La columna derecha de D2/D3 superpone los contornos $P(y=1|x) = 0.5$
solo para runs con acc $\geq 0.95$, sobre una nube de puntos de todos los
datasets.

\subsection{Experimento E: Robustez de $\nu^*$ (make\_moons)}

\textbf{Objetivo.} Verificar que la distribución óptima de parámetros $\nu^*$
es robusta a distintas condiciones de entrenamiento: distintas inicializaciones
y distintos datasets deben producir distribuciones $\nu^*$ similares.

\textbf{Protocolo.} Dos sub-experimentos con $n_\text{seeds} = 10$,
$\varepsilon = 0.01$, 500 épocas y \textbf{pSGLD} (\texttt{use\_sgld=True},
para explorar la distribución de Gibbs $\nu^*$ en lugar del MAP):
\begin{itemize}
  \item \textbf{E-1:} \texttt{data\_seed}$=42$ fijo, \texttt{init\_seed}
        $\in \{0,\ldots,9\}$ variable.
  \item \textbf{E-2:} \texttt{init\_seed}$=2$ fijo, \texttt{data\_seed}
        $\in \{0,\ldots,9\}$ variable.
\end{itemize}

\textbf{Magnitud visualizada.} Para cada run se extraen los pesos de salida
$a_0^m \in \R^2$ de las $M=64$ neuronas y se calcula la \emph{importancia}:
\begin{equation}
  I^m = \norm{a_0^m}_2.
\end{equation}
Las importancias se ordenan de mayor a menor y se representan para cada seed,
junto con la media $\pm 1\sigma$ sobre las 10 seeds (figura 1×2):
\begin{itemize}
  \item \textbf{Panel izquierdo (E-1):} importancias con init variable,
        datos fijos. Banda estrecha $\Rightarrow$ $\nu^*$ es robusta a la
        inicialización.
  \item \textbf{Panel derecho (E-2):} importancias con datos variables,
        init fija. Banda estrecha $\Rightarrow$ $\nu^*$ es una propiedad
        de la geometría del problema, no del dataset concreto.
\end{itemize}

Si el Meta-Teorema~1 se cumple (minimizador único estable), las curvas de
importancia deben ser cualitativamente similares entre seeds: mismo ``codo''
pronunciado en los primeros ranks y misma caída.

\subsection{Experimento F: Simetría de $\nu^*$ en \texttt{make\_circles}}

\textbf{Objetivo.} Verificar que en \texttt{make\_circles} (dataset con
simetría rotacional $\text{SO}(2)$) la convergencia del modelo es robusta
tanto a variaciones del dataset como de la inicialización.

\textbf{Protocolo.} Dos sub-experimentos con $n_\text{seeds} = 10$,
$n_\text{epochs} = 700$, $\varepsilon = 0.01$ y \textbf{pSGLD}
(\texttt{use\_sgld=True}):
\begin{itemize}
  \item \textbf{F1:} \texttt{init\_seed}$= 4$ fija, \texttt{data\_seed}
        variable.
  \item \textbf{F2:} \texttt{data\_seed}$= 42$ fijo, \texttt{init\_seed}
        variable.
\end{itemize}

\textbf{Figura.} Se muestra un único panel con las curvas de convergencia
$J \pm 1\sigma$ para F1 y F2 superpuestas. La banda más estrecha de F1 vs.\
F2 (o viceversa) indica cuál fuente de variabilidad domina la robustez del
entrenamiento.

\textbf{Métrica de convergencia.} Como pSGLD explora la distribución de Gibbs
$\nu^* \propto \exp(-J/\varepsilon)$ y no converge puntualmente a $J^*$, el
nivel de pérdida se reporta como $\bar{J}_\text{final}$, la media de las
últimas 50 épocas, que estima el nivel estacionario de la cadena de Markov.

\subsection{Experimento G: Problema de convergencia ($N \to \infty$)}

\textbf{Objetivo.} Verificar empíricamente que el valor óptimo del problema
de campo medio finito $J^*_N$ converge al valor óptimo del problema de campo
medio con infinitas partículas $J^*_\infty$ cuando $N \to \infty$, y estimar
la tasa de convergencia.

\textbf{Marco teórico.} El paper \cite{daudin2025} establece que
\begin{equation}
  J^*_N \;\xrightarrow{N\to\infty}\; J^*_\infty.
\end{equation}
Una tasa natural esperada por analogía con la ley de los grandes números es
$\text{gap}(N) := J^*_\infty - J^*_N \sim C \cdot N^{-\alpha}$ con
$\alpha \approx 0.5$. El experimento busca la tasa empírica.

\textbf{Protocolo.} Se consideran $N \in \{25, 50, 100, 200, 400, 800\}$
partículas y $\varepsilon \in \{0, 0.01\}$. Para cada combinación $(N, \varepsilon)$
se realizan 20 entrenamientos independientes (\texttt{n\_seeds}$=20$) con
700 épocas y \textbf{pSGLD} (\texttt{use\_sgld=True}). Se usa el optimizador
pSGLD (no Adam) para que el nivel de ruido sea coherente con el parámetro
$\varepsilon$: con $\varepsilon=0$ el ruido es exactamente cero.

\textbf{Estimador de $J^*_\infty$.} El valor $J^*_\infty$ es inaccesible
directamente; se aproxima mediante un \emph{test set oráculo}: se evalúa la
pérdida BCE (sin término de regularización) del modelo entrenado sobre un
conjunto de test de $N_\text{test} = 4000$ puntos i.i.d.\ de make\_moons.
La hipótesis es que, para $N$ grande, el BCE en test tiende a $J^*_\infty$
(la pérdida de Bayes en la distribución de datos). El gap se define como:
\begin{equation}
  \text{gap}(N) := J^*_\infty - J^*_N \;\approx\;
  \text{BCE}_\text{test}(N_\text{large}) - \text{BCE}_\text{test}(N).
\end{equation}
En la práctica, $J^*_\infty$ se aproxima con la media de BCE\_test sobre las
20 semillas para $N=800$.

\textbf{Figura (G\_convergence\_problem.png).} Layout 1×2:
\begin{itemize}
  \item \textbf{G1 — Convergencia de la pérdida en test:} BCE\_test (media
        $\pm 1\sigma$ sobre 20 seeds) vs.\ $N$ en escala log-log, para
        $\varepsilon=0$ y $\varepsilon=0.01$. Evidencia visual de que
        BCE\_test$\,(N) \to J^*_\infty$ al crecer $N$.
  \item \textbf{G2 — Tasa de convergencia del gap:} $\log\,\text{gap}(N)$ vs.\
        $\log N$ con ajuste lineal (mínimos cuadrados sobre $N \geq 100$),
        cuya pendiente es $-\alpha$. Los resultados empíricos son:
        $\varepsilon=0$: $\alpha \approx 0.78$;
        $\varepsilon=0.01$: $\alpha \approx 0.85$.
        Ambas tasas son más rápidas que la referencia $N^{-0.5}$.
\end{itemize}

\textbf{Checkpoint incremental.} Para soportar interrupciones (sesiones de
Colab), los resultados se guardan en \texttt{figuras/G\_results.npz} tras
completar cada valor de $N$. Al relanzar, las combinaciones $(N,\varepsilon)$
ya calculadas se omiten automáticamente.

% ============================================================
\section{Resumen de la arquitectura}
\label{sec:resumen}

\begin{center}
\begin{tabular}{ll}
\toprule
Bloque & Descripción \\
\midrule
Datasets & make\_moons (A–E, G) y make\_circles (F), $N=400$, normalizados \\
ODE & $dX_t/dt = F(X_t,t)$, $t \in [0,1]$, 10 pasos RK4 \\
Campo $F$ & $(1/M)\sum_m \tanh(a_1^m \cdot x + a_2^m t + b_1^m)\,a_0^m$,\; $M=64$ \\
Parámetros ODE & $W_1 \in \R^{64\times3}$, $b_1\in\R^{64}$, $W_0\in\R^{2\times64}$ (384 total)\\
Clasificador & $\ell = W \cdot X_T + b$,\; $W \in \R^{1\times 2}$,\; $b \in \R$ (3 total) \\
Total parámetros & 387 \\
Pérdida & $J = \text{BCE} + \varepsilon \cdot \mathcal{E}$,\; $\varepsilon = 0.01$ \\
Optimizador (A, D) & Adam \cite{kingma2015adam},\; $\alpha_0=0.01$,\; cosine annealing \\
Optimizador (C) & SGD,\; $\alpha=0.01$ constante,\; sin ruido \\
Optimizador (B, E, F, G) & pSGLD \cite{li2016psgld},\; $\theta \leftarrow \text{Adam}(\theta,\nabla J) + \sqrt{2\eta\varepsilon M}\,\xi$,\; $M = \min(1/\sqrt{\hat{v}}+\delta, 1)$ \\
Gradient clip & $G_{\max}=5$ (todos los modos) \\
\bottomrule
\end{tabular}
\end{center}

% ============================================================
\begin{thebibliography}{9}

\bibitem{daudin2025}
  S.~Daudin and F.~Delarue,
  \textit{Genericity of the Polyak-{\L}ojasiewicz inequality for mean-field
  Neural ODEs with entropic regularization},
  arXiv:2507.08486, 2025.

\bibitem{welling2011sgld}
  M.~Welling and Y.~W.~Teh,
  \textit{Bayesian Learning via Stochastic Gradient Langevin Dynamics},
  in \textit{International Conference on Machine Learning (ICML)}, 2011.

\bibitem{li2016psgld}
  C.~Li, C.~Chen, D.~Carlson, and L.~Carin,
  \textit{Preconditioned Stochastic Gradient Langevin Dynamics for Deep
  Neural Networks},
  in \textit{AAAI Conference on Artificial Intelligence}, 2016.
  \url{https://arxiv.org/abs/1512.07666}

\bibitem{kingma2015adam}
  D.~P.~Kingma and J.~Ba,
  \textit{Adam: A Method for Stochastic Optimization},
  in \textit{International Conference on Learning Representations (ICLR)},
  2015.
  \url{https://arxiv.org/abs/1412.6980}

\bibitem{loshchilov2017sgdr}
  I.~Loshchilov and F.~Hutter,
  \textit{SGDR: Stochastic Gradient Descent with Warm Restarts},
  in \textit{International Conference on Learning Representations (ICLR)},
  2017.
  \url{https://arxiv.org/abs/1608.03983}

\bibitem{chen2018node}
  R.~T.~Q.~Chen, Y.~Rubanova, J.~Bettencourt, and D.~Duvenaud,
  \textit{Neural Ordinary Differential Equations},
  in \textit{Advances in Neural Information Processing Systems (NeurIPS)},
  2018.
  \url{https://arxiv.org/abs/1806.07366}

\bibitem{polyak1963}
  B.~T.~Polyak,
  \textit{Gradient methods for the minimisation of functionals},
  \textit{USSR Computational Mathematics and Mathematical Physics},
  3(4):864--878, 1963.

\bibitem{lojasiewicz1963}
  S.~Łojasiewicz,
  \textit{Une propriété topologique des sous-ensembles analytiques réels},
  in \textit{Colloques Internationaux du CNRS}, 117:87--89, 1963.

\bibitem{villani2009}
  C.~Villani,
  \textit{Optimal Transport: Old and New},
  Springer, 2009.

\end{thebibliography}

\end{document}
